% Preamble — Berkeley/NCAR, matching the house style of
% ~/DavidFillmore/Media/decks/agentic_ai (which in turn follows
% ~/EarthSystem/CheMPAS-Slides/ceres_talk). Trimmed for this deck: no
% beamerposter, no CLI transcript panels, no command badges.
\documentclass[aspectratio=169]{beamer}
% width=2.6cm (Berkeley default 2cm): the section names "AOD Constraints" and
% "Band Properties" wrap at the default width.
\usetheme[width=2.6cm]{Berkeley}
\useinnertheme{default}

\usepackage[T1]{fontenc}
\usepackage{lmodern}

% Tectonic reads UTF-8 natively, so a literal em dash has no mapping into the
% T1-encoded body font and is silently dropped ("Missing character"). inputenc
% is a no-op on a unicode-native engine, so map it by catcode. Slide sources
% here use `---` throughout; this is insurance against a pasted em dash.
\catcode`\—=\active
\def—{\textemdash}

\usepackage{graphicx}
\usepackage{tikz}
\usetikzlibrary{shapes.geometric, arrows.meta, positioning, fit, backgrounds,
                calc, decorations.pathreplacing, patterns}
\usepackage{booktabs}
\usepackage{amsmath}
\usepackage{amssymb}

% Palette, symbolic seam, logo binding and graphics path.
% NASA theme seam — palette, logo binding, graphics path.
% After CERES-MIP/slides/colors.tex, but the whole deck is NASA: the seam
% below binds every slide's chrome to NASA blue (what \nasaslide did for a
% single frame in Media/decks/agentic_ai). NCAR inks stay defined for use
% inside diagrams (model-side boxes) only.

% --- NASA palette --------------------------------------------------------
\definecolor{nasaBlue}{HTML}{0B3D91}
\definecolor{nasaRed}{HTML}{FC3D21}

% --- NSF NCAR brand palette (diagram inks only; not deck chrome) ---------
\definecolor{ncarSpace}{HTML}{011837}
\definecolor{ncarDarkBlue}{HTML}{00357A}
\definecolor{ncarBlue}{HTML}{0A5DDA}
\definecolor{ncarAqua}{HTML}{00A2B4}
\definecolor{ncarLightBlue}{HTML}{CEDFF8}
\definecolor{ncarLightGray}{HTML}{F1F0EE}
\definecolor{ncarOrange}{HTML}{FF8C00}
\definecolor{ncarYellow}{HTML}{FFDD31}

% --- Symbolic seam: slide patterns use ONLY these (portability contract) ---
\colorlet{accent}{nasaBlue}            % structure, frametitle, bullets
\colorlet{accentSoft}{nasaBlue!40!white}% secondary tier
\colorlet{accentDeep}{nasaBlue!60!black}% deep tier
\colorlet{sidebarBg}{nasaBlue!35!black} % Berkeley sidebar ground
\colorlet{bannerBg}{nasaBlue}           % headline banner boxes

% --- Logo binding (bare filename; resolved via \graphicspath below) ---
\newcommand{\themelogo}{nasa_logo.pdf}

% --- Font binding ---
% helvet: Poppins is not shipped with Tectonic.
\usepackage{helvet}
\renewcommand{\familydefault}{\sfdefault}

\graphicspath{{graphics/}}


% --- Beamer chrome -------------------------------------------------------
% Written against the symbolic seam (accent / accentSoft / accentDeep /
% sidebarBg) rather than the literal ncar* names, so a future per-slide
% recolor can redirect one frame: beamer stores the color NAME and resolves
% it at typeset time, which a group-local \colorlet redirects.
\setbeamercolor{structure}{fg=accent}
\setbeamercolor{palette primary}{bg=accent, fg=white}
\setbeamercolor{palette secondary}{bg=accentSoft, fg=white}
\setbeamercolor{palette tertiary}{bg=accentDeep, fg=white}
\setbeamercolor{palette quaternary}{bg=sidebarBg, fg=white}
\setbeamercolor{frametitle}{bg=accent, fg=white}
\setbeamercolor{sidebar}{bg=sidebarBg}
\setbeamercolor{section in sidebar}{fg=accentSoft!30!white}
\setbeamercolor{section in sidebar shaded}{fg=accentSoft}
\setbeamerfont{section in sidebar}{series=\bfseries, size*={7.5pt}{8.4pt}}
\setbeamerfont{section in sidebar shaded}{series=\mdseries, size*={7.5pt}{8.4pt}}
\setbeamercolor{author in sidebar}{fg=white}
\setbeamercolor{title in sidebar}{fg=white}
% Sidebar = section nav only, starting at the top. Berkeley's stock template
% stacks short title + short author + ~4em of glue above the nav; this
% replacement drops all of it (emptying the short forms is not enough — the
% glue stays).
\makeatletter
\setbeamertemplate{sidebar left}{%
  \vskip0.5em%
  \insertverticalnavigation{\beamer@sidebarwidth}%
  \vfill
}
\makeatother
\setbeamercolor{section in toc}{fg=accentSoft}
\setbeamercolor{block title}{bg=accent, fg=white}
\setbeamercolor{block body}{bg=ncarLightGray, fg=ncarSpace}
\setbeamercolor{block title alerted}{bg=ncarOrange, fg=white}
\setbeamercolor{block title example}{bg=accentSoft, fg=white}
\setbeamercolor{item}{fg=accent}
\setbeamercolor{subitem}{fg=accentSoft}
\setbeamercolor{subsubitem}{fg=accentDeep}

\setbeamertemplate{navigation symbols}{}

% \fitline{<content>} — set <content> on one line within \linewidth, shrinking
% only if it would otherwise wrap. Never enlarges. Every diagram on this deck
% is wrapped in it, so the tikz coordinate extent is authored just under the
% body column and the font sizes written there are what reach the slide.
% No blank lines inside it: it measures with \settowidth, and a \par there
% ends the measurement mid-box.
\newlength{\fitlinewd}
\newcommand{\fitline}[1]{%
  \settowidth{\fitlinewd}{#1}%
  \ifdim\fitlinewd>\linewidth
    \resizebox{\linewidth}{!}{#1}%
  \else
    #1%
  \fi
}

% \fitframe{<diagram>} — the workhorse every content frame uses.
%
% \fitline alone is not enough here. It constrains WIDTH only, so a diagram
% authored just inside the column but too tall still overruns the page, and
% beamer reports it as "Overfull \vbox" with the bottom of the slide silently
% cropped. Each of these frames pairs a wide diagram with a wrapped text strip
% whose height depends on how the prose happens to break — which is exactly
% the kind of thing that shifts under a one-word edit.
%
% adjustbox's max size scales proportionally to fit BOTH bounds and leaves
% anything already smaller untouched, so a frame can be edited freely and
% still lands on the page. \bodyheight is the room left under the frametitle
% bar.
% NOTE ON TYPE SIZE. max size only ever shrinks, so any diagram authored
% larger than the body box gets scaled DOWN and every font in it shrinks with
% it — a \scriptsize label inside a picture scaled to 0.8 lands near 6 pt and
% is unreadable projected. So frames are authored to fit at 1:1 (about 14 cm
% wide, \bodyheight tall) and their labels start at \normalsize, never
% smaller. If a frame needs shrinking to fit, cut content instead.
%
% \bodyheight is 6.50cm. The body region under the frametitle measures ~7.43cm
% on paper but only ~6.5cm is actually usable — 7.00cm overran the page by
% 0.47cm. Measured across the deck, the smallest achievable gap between the
% frametitle bar and the first ink is ~0.49cm: that residue is beamer's own
% frametitle padding, not slack in the picture, so do not chase it. Frames
% measuring 0.9cm+ DO have reclaimable space. Beamer centres the picture in
% what is left, so the balance is
% frametitle: beamer centres the picture in that region, so the balance is
% margin. It was 6.15cm, which was shrinking the three band frames (6.4 units
% tall) to 0.96 and taking every font in them down with it.
\usepackage{adjustbox}
\newlength{\bodyheight}
\setlength{\bodyheight}{6.50cm}
\newcommand{\fitframe}[1]{\adjustbox{max size={\linewidth}{\bodyheight}}{#1}}

% --- Molecular formulae --------------------------------------------------
% \mgap is the gap inside a chemical formula. ONE rule, applied everywhere:
% split a formula at the bonds BETWEEN chemically distinct moieties, and keep
% each functional group tight.
%
%   split      Cl\mgap O\mgap NO$_2$    chlorine - oxygen - nitro
%              H\mgap NO$_3$            proton - nitrate
%              H\mgap CN                proton - cyanide
%   tight      O$_3$  O$_2$  H$_2$O  CO$_2$  CO  CH$_4$  N$_2$O  NO$_2$
%              N$_2$O$_5$
%
% Run together, ClONO$_2$ is genuinely hard to read when projected — "ClO"
% scans as C-l-O — and the same applies to HNO$_3$ and HCN. N$_2$O$_5$ stays
% tight because its conventional formula is not a concatenation of a distinct
% atom and a group; splitting it would need the structural form O$_2$N-O-NO$_2$,
% which is not what an audience expects to see. CFC-11 and CFC-12 are names,
% not formulae.
%
% Routed through a macro so the convention is enforced in one place: widen or
% narrow the gap here, not in fourteen call sites.
\newcommand{\mgap}{\,}

% \codepath{...} — typewriter formatting for path and dataset-version refs.
\newcommand{\codepath}[1]{\texttt{\small #1}}

% --- Link convention (deck-wide) -----------------------------------------
% EVERY link renders as REGULAR TEXT, UNDERLINED. No color shift, no italics,
% no beamer link box — the rule is the only affordance. Routed through a macro
% so the convention is enforced by construction.
\newcommand{\decklink}[2]{\href{#1}{\mbox{\underline{#2}}}}
\hypersetup{colorlinks=false, pdfborder={0 0 0}}


% \regpanel{<pdf>} — one regional daily time-series panel (656x296pt DAVINCI
% plots, ~7pt margins so no trim), 0.48 of the column width so a 2x2 fits
% under \bodyheight. Used by the per-region 2008 Test Run frames.
\newcommand{\regpanel}[1]{\includegraphics[width=0.48\linewidth]{#1}}

% \zonalpanel{<pdf>} — one annual zonal-mean latitude-pressure section
% (656x368pt DAVINCI plots, ~7pt margins so no trim). Height-limited to
% 3.1cm like \compmap (frames/07): at 0.48 of the column width these would
% stand 3.4cm and two rows would overrun \bodyheight. Used by the 2008 Test
% Run zonal frames (12b, 12c).
\newcommand{\zonalpanel}[1]{\includegraphics[height=3.1cm]{#1}}

% --- Bullet slides ---------------------------------------------------------
% The intro slides (02-05) are beamer blocks of bullets, two columns. Stock
% block padding plus \small wrapped bullets overran \bodyheight by 2-3cm, so
% blocks are trimmed top and bottom and items sit closer together.
% The trims go OUTSIDE the box (before block begin, after block end): put
% inside, the negative space pulls the grey background up over the last
% line of the body and the text spills below the box.
\addtobeamertemplate{block begin}{\vspace{-0.3em}}{}
\addtobeamertemplate{block end}{}{\vspace{-0.6em}}
\setlength{\leftmargini}{1.1em}
\setbeamertemplate{itemize/enumerate body begin}{\setlength{\itemsep}{0.05em}}
